Use barycentric weights $\alpha,\beta,\gamma>0$ with sum $1$, so $P=\alpha A+\beta B+\gamma C$. Along $BE$, the ratio $\dfrac{BP}{PE}=\dfrac{1-\beta}{\beta}$, hence $\beta=\dfrac{2}{7}$. Similarly $\gamma=\dfrac{3}{10}$. Therefore $\alpha=1-\dfrac{2}{7}-\dfrac{3}{10}=\dfrac{29}{70}$ and $\dfrac{AP}{PD}=\dfrac{1-\alpha}{\alpha}=\dfrac{41}{29}$. Thus $\boxed{p+q=70}$.
\begin{asy}
texpreamble("\usepackage{amsmath,amssymb}");
size(240);
defaultpen(fontsize(9pt));
pair A=(0,5),B=(0,0),C=(7,0),P=29/70.0*A+2/7.0*B+3/10.0*C,D=(2/7.0*B+3/10.0*C)/(1-29/70.0),E=(29/70.0*A+3/10.0*C)/(1-2/7.0),F=(29/70.0*A+2/7.0*B)/(1-3/10.0);draw(A--B--C--cycle);draw(A--D);draw(B--E);draw(C--F);label("A",A,N);label("B",B,SW);label("C",C,SE);label("D",D,S);label("E",E,NE);label("F",F,W);label("P",P,SE);
\end{asy}
